\chapter[Statements and mathematical structures]{From ordinary statements to mathematical structures}
A mathematical proof is a sequence of statements whose conclusion follows
from explicitly declared assumptions. Start with the sentence
\emph{If a spike occurred in an interval, the interval contains a spike}.
The implication is meaningful only after the interval and the membership
relation have been defined.

\begin{definition}
A \emph{relation} from $A$ to $B$ is a subset $R\subseteq A\times B$.
It is a \emph{function} $f:A\to B$ when every $a\in A$ occurs as
the first coordinate of exactly one pair $(a,b)\in R$.
\end{definition}
For example, a spike train of a unit is a subset of its observation
time interval. A classification of spike waveforms by cluster labels is
a function only when every input waveform receives exactly one label.
This does not imply that the label identifies a physiological neuron.

\entry{1}{De Morgan's laws}
{For propositions $P,Q$, $\neg(P\land Q)\Leftrightarrow
(\neg P)\lor(\neg Q)$; for sets,
$(A\cup B)^c=A^c\cap B^c$, relative to an ambient set.}
{Transform selection criteria and complements of detected event sets.}
\entry{2}{Quantifier duality}
{For any predicate $P$, $\neg(\forall x\,P(x))\Leftrightarrow
\exists x\,\neg P(x)$; likewise
$\neg(\exists x\,P(x))\Leftrightarrow\forall x\,\neg P(x)$.}
{Express universal reliability claims and their counterexamples.}
\entry{3}{Mathematical induction (principle)}
{If $P(0)$ holds and $P(n)\Rightarrow P(n+1)$ for $n\in\N$,
then $P(n)$ holds for all $n\in\N$.}
{Prove recurrences, finite sums, binning bounds, and iterative invariants.}
\entry{4}{Recursion theorem on the natural numbers}
{With a starting value and a step map, there is a unique
sequence satisfying the corresponding recursive definition
in the standard set-theoretic setting.}
{Construct discrete filters, numerical solvers, and event counters.}
\entry{5}{Equivalence relations and partitions}
{Reflexive, symmetric, transitive relations correspond to partitions
into disjoint nonempty equivalence classes.}
{Distinguish clustering as a partition from physiological identification.}
\entry{6}{Well-ordering principle}
{Every nonempty subset of $\N$ has a least element.}
{Prove termination of monotone integer-index algorithms.}
\entry{7}{Cantor--Schroeder--Bernstein theorem}
{If injections $A\to B$ and $B\to A$ exist, then a bijection
$A\to B$ exists.}
{Understand cardinality and information-preserving relabelings.}
\entry{8}{Cantor's diagonal argument}
{There is no surjection $\N\to(0,1)$; the real interval
is uncountable.}
{Understand the distinction between finite sampled arrays and continuum models.}
\entry{9}{Least-upper-bound property (completeness axiom for $\R$)}
{Every nonempty subset of $\R$ bounded above has a real supremum.}
{Build limits, continuity, completeness, and convergence of numerical approximations.}
\entry{10}{Archimedean property and density of $\Q$}
{For real $x$ there is an integer $n>x$; between any
$a<b$ in $\R$ lies some $q\in\Q$.}
{Approximate continuous times and amplitudes by rational values, without confusing
approximation with exact sampling.}

\begin{example}[Equivalence classes of spike timestamps]
For a fixed bin width $\Delta>0$, let $s\sim t$ when
$\lfloor s/\Delta\rfloor=\lfloor t/\Delta\rfloor$.
This relation is reflexive, symmetric, and transitive.
Each equivalence class represents one temporal bin, not one spike.
A spike-count sequence counts members of each class.
\end{example}

\begin{caution}
Membership ($x\in A$), inclusion ($A\subseteq B$), and association
under a binary relation ($xRy$) are distinct concepts. An observed
cross-correlation is not evidence that a causal or anatomical edge exists.
\end{caution}
% --- expanded edition (October 2026) ---
\section{Proofs of selected foundational results}
The catalogue entries above state results compactly. This section proves
those whose arguments are short enough to be reproduced by hand and
whose structure reappears later in the book.

\begin{proof}[Proof of Result 1 for sets]
Let $U$ be the ambient set and $x\in U$. Then $x\in(A\cup B)^c$
iff $\neg(x\in A\lor x\in B)$ iff $(x\notin A)\land(x\notin B)$
by the propositional De Morgan law, iff $x\in A^c\cap B^c$.
Since the two sets have the same members they are equal. The
propositional law itself is verified by its four-row truth table.
\end{proof}

\begin{proposition}[Result 6 implies Result 3]
The well-ordering principle implies the principle of induction.
\end{proposition}
\begin{proof}
Assume $P(0)$ and $P(n)\Rightarrow P(n+1)$ for all $n$, and suppose
the set $F=\{n\in\N:\neg P(n)\}$ is nonempty. Let $m$ be its least
element. Then $m\ne0$ because $P(0)$ holds, so $m-1\in\N$ and
$m-1\notin F$, i.e.\ $P(m-1)$ holds. The induction step gives $P(m)$,
contradicting $m\in F$.
\end{proof}

\begin{proof}[Proof of Result 8]
Suppose $g:\N\to(0,1)$ is surjective and write each $g(n)$ by its
decimal expansion $0.d_{n1}d_{n2}\ldots$, choosing the expansion that
does not end in repeated $9$'s. Define $y=0.e_1e_2\ldots$ with
$e_n=5$ if $d_{nn}\ne5$ and $e_n=4$ if $d_{nn}=5$. Then $y\in(0,1)$,
its expansion is of the chosen kind, and it differs from $g(n)$ in the
$n$th digit for every $n$. Hence $y$ is not in the image of $g$.
\end{proof}

\begin{proof}[Proof sketch of Result 7]
Let $f:A\to B$ and $g:B\to A$ be injective. Call $a\in A$ an
\emph{$A$-stopper} if, tracing preimages
$a,\;g^{-1}(a),\;f^{-1}(g^{-1}(a)),\ldots$ as far as they exist,
the chain terminates in $A\setminus g(B)$. Define $h(a)=f(a)$ for
$A$-stoppers and $h(a)=g^{-1}(a)$ otherwise. One checks that every
non-stopper lies in $g(B)$, that $h$ is injective on each of the two
pieces, and that the two images partition $B$. Thus $h$ is a bijection.
\end{proof}
Result 7 is used implicitly whenever one argues that two index sets
``have the same size'' by exhibiting encodings in each direction, for
example that finite spike rasters over a countable set of bins can be
enumerated.

\section{Relations as data structures}
Most data objects in an electrophysiology pipeline are relations or
functions whose properties should be checked rather than assumed.
\begin{itemize}[leftmargin=2em]
\item A spike table with columns (unit, time) is a relation
$R\subseteq U\times\R$. It is the graph of a function $U\to\R$ only
if each unit fires exactly once.
\item A channel map is a function from channel indices to physical
positions. It must be injective (no two channels at the same place)
for current-source density estimates to make sense.
\item A trial table is a function from trial indices to
(condition, onset) pairs. If it is not injective in onset, two trials
share a time window and their data are not independent.
\item A cluster assignment is a function from detected waveforms to
labels, equivalently a partition (Result 5). A \emph{merge} of two
clusters replaces the partition by a coarser one; a \emph{split}
replaces it by a finer one.
\end{itemize}

\begin{proposition}[Agreement between two sortings]
Let $\mathcal P$ and $\mathcal Q$ be two partitions of the same finite
set of detected spikes $S$. The meet
$\mathcal P\wedge\mathcal Q=\{P\cap Q:P\in\mathcal P,\ Q\in\mathcal Q,\ P\cap Q\ne\varnothing\}$
is a partition of $S$ that refines both.
\end{proposition}
\begin{proof}
Each $s\in S$ lies in exactly one $P$ and one $Q$, hence in exactly one
$P\cap Q$. Blocks are nonempty by construction and pairwise disjoint
because distinct pairs $(P,Q)$ differ in at least one coordinate, whose
blocks are disjoint. Each $P\cap Q$ is contained in $P$ and in $Q$.
\end{proof}
The confusion matrix $M_{PQ}=|P\cap Q|$ used to compare sorters is
the table of block sizes of the meet. Agreement indices such as the
adjusted Rand index are functions of this table and answer the
question ``how similar are these two partitions?'', not ``which
partition corresponds to neurons?''.

\begin{example}[Bins as equivalence classes, continued]
For the binning relation $s\sim t\iff\lfloor s/\Delta\rfloor=\lfloor t/\Delta\rfloor$,
the quotient map $t\mapsto\lfloor t/\Delta\rfloor$ is a function from
times to integers. It is not injective, so binning cannot be undone.
Two spikes of different units in the same bin are \emph{coincident at
resolution $\Delta$}; the same pair may be non-coincident at resolution
$\Delta/2$. Coincidence is therefore a relation that depends on a
parameter, and every synchrony claim must state that parameter.
\end{example}

\section*{Exercises for this chapter}
\begin{exercise}
Prove the set form of De Morgan's law for an arbitrary indexed family
$(A_i)_{i\in I}$.
\end{exercise}
\begin{exercise}
Show that the induction principle implies the well-ordering principle,
completing the equivalence of Results 3 and 6.
\end{exercise}
\begin{exercise}
Prove that the set of finite subsets of $\N$ is countable, and that
the set of all subsets of $\N$ is not.
\end{exercise}
\begin{exercise}
Two sorters produce partitions of 10{,}000 spikes into 12 and 15
clusters. What is the largest and the smallest possible number of
blocks in their meet?
\end{exercise}
